
\section{formula di Cauchy-Binet}
\begin{theorem}[formula di Cauchy-Binet]
  Siano $A$ una matrice $m\times n$ e $B$ una matrice $n\times m$.
  Allora 
  \[
    \det (AB) = \sum_{\vec k} \det A^{\vec k}\cdot \det B_{\vec k}
  \]
  dove 
  $\vec k = (k_1,\dots,k_m) \in \ENCLOSE{1,\dots,n}^m$ 
  è un elenco crescente di indici: $1\le k_1 < \dots < k_m\le n$
  e $A^{\vec k}$ e $B_{\vec k}$ sono le matrici $m\times m$ che 
  rappresentano i \emph{minori} di $A$ e $B$:
  \[
    (A^{\vec k})_{i,j} = A_{i,k_j}, \qquad
    (B_{\vec k})_{i,j} = B_{k_i, j}.
  \]
\end{theorem}
%
\begin{proof}
Denotiamo con $\Sigma_m$ il gruppo delle permutazioni 
di $\ENCLOSE{1,\dots,m}$ 
e con $(-1)^\sigma$ il segno della permutazione.
Si ha
\begin{align*}
  \det (AB) &= \sum_{\sigma \in \Sigma_m} (-1)^\sigma \prod_{i=1}^m (AB)_{i,\sigma(i)} \\
    &= \sum_\sigma (-1)^\sigma \prod_{i=1}^m \sum_{k=1}^n A_{i,k} B_{k,\sigma(i)} \\
    &= \sum_\sigma (-1)^\sigma \sum_{k_1=1}^n \dots \sum_{k_m=1}^n \prod_{i=1}^m A_{i,k_i} B_{k_i,\sigma(i)} \\
    &= \sum_{k_1,\dots,k_m=1}^n \prod_{i=1}^m A_{i,k_i} \sum_\sigma (-1)^\sigma \prod_{i=1}^m B_{k_i,\sigma(i)} \\
    &= \sum_{\vec k\in \ENCLOSE{1,\dots,n}^m} \prod_{i=1}^m A_{i,k_i} \det B_{\vec k}
\end{align*}
Notiamo ora che se $\vec k = (k_1,\dots,k_m)$ non è iniettiva, cioè se $k_i = k_j$ per qualche 
$i\neq j$, la matrice $B_{\vec k}$ ha due righe uguali e quindi $\det B_{\vec k}=0$. 
Possiamo dunque restringere $\vec k$ alle funzioni iniettive.
Se $\vec k\colon\ENCLOSE{1,\dots,m}\to\ENCLOSE{1,\dots,n}$ è iniettiva c'è una 
unica permutazione $\phi$ tale che $i\mapsto k_{\phi(i)}$ è strettamente crescente.
Dunque ogni $\vec k$ iniettiva si ottiene prendendo tutte le $\vec k$ strettamente crescenti 
e componendole con tutte le permutazioni $\phi$ di $\ENCLOSE{1,\dots,m}$. 
Dunque
\begin{align*}
  \det(AB) 
    &= \sum_{1\le k_1<\dots<k_m\le n}\sum_{\phi\in \Sigma_m}
    \prod_i A_{i,k_{\phi(i)}} \det B_{\vec k\circ \phi}\\
    &= \sum_{\phi\in \Sigma_m} \sum_{1\le k_1<\dots<k_m\le n}
    \prod_i A_{i,k_{\phi(i)}} (-1)^\phi \det B_{\vec k}\\
    &= \sum_{1\le k_1<\dots<k_m\le n} 
    \det A^{\vec k} \det B_{\vec k}
\end{align*}
avendo usato il fatto che $\det B_{\vec k\circ \phi}$
si ottiene da $\det B_{\vec k}$ scambiando le righe 
secondo la permutazione $\phi$, ogni scambio di righe 
cambia segno al determinante e dunque si ha $\det B_{\vec k\circ \phi} = (-1)^\phi \det B_{\vec k}$.
\end{proof}


